% Compile with achemso, for example:
%   pdflatex jctc_unique_spin_vbscf
%   pdflatex jctc_unique_spin_vbscf
%
% This draft uses the historical-but-stable ACS `achemso` workflow for local
% drafting. Update journal metadata / author block as needed before submission.

\documentclass[journal=jctcce,manuscript=article]{achemso}

\usepackage[version=3]{mhchem}
\usepackage{amsmath,amssymb,amsfonts,mathtools,bm}
\usepackage{graphicx}
\usepackage{booktabs}
\usepackage{multirow}
\usepackage{array}
\usepackage{url}
\usepackage{xcolor}
\usepackage{siunitx}

\DeclareMathOperator{\Tr}{Tr}
\DeclareMathOperator{\diag}{diag}
\DeclareMathOperator{\rank}{rank}
\DeclareMathOperator{\adj}{adj}
\newcommand{\Frob}[2]{\left\langle #1, #2 \right\rangle_{\mathrm F}}
\newcommand{\na}{N_{\alpha}}
\newcommand{\nb}{N_{\beta}}
\newcommand{\ndet}{N_{\mathrm{det}}}
\newcommand{\nstr}{N_{\mathrm{str}}}
\newcommand{\nstate}{N_{\mathrm{state}}}
\newcommand{\npair}{N_{\mathrm{pair}}}
\newcommand{\Cmat}{\mathbf{C}}
\newcommand{\Sa}{\mathbf{S}^{\alpha}}
\newcommand{\Sb}{\mathbf{S}^{\beta}}
\newcommand{\Ha}{\mathbf{H}^{\alpha}}
\newcommand{\Hb}{\mathbf{H}^{\beta}}
\newcommand{\Honea}{\mathbf{H}^{1e,\alpha}}
\newcommand{\Honeb}{\mathbf{H}^{1e,\beta}}
\newcommand{\Ua}[1]{\mathbf{U}^{\alpha}_{#1}}
\newcommand{\Ub}[1]{\widetilde{\mathbf{U}}^{\beta}_{#1}}
\newcommand{\wstate}{w_n}
\newcommand{\En}{E_n}
\newcommand{\todo}[1]{\textcolor{red}{[TODO: #1]}}

\author{Tao Xia}
\author{Wei Wu}
\author{Chen Zhou}
\email{TODO-corresponding-author@xmu.edu.cn}
\affiliation{Department of Chemistry, Xiamen University, Xiamen 361005, Fujian, China}

\title{An Exact Unique-Spin-String Reformulation of Determinant Algebra in Valence-Bond Self-Consistent Field Theory}

\begin{abstract}
Valence-bond self-consistent field (VBSCF) theory provides a chemically
transparent description of electronic structure, but the practical evaluation
of overlap, Hamiltonian, and analytic derivative terms is often limited by
repeated summations over full determinant pairs. In typical VBSCF expansions,
however, the determinant list has strong internal redundancy: many full
determinants share the same $\alpha$-spin string, the same $\beta$-spin
string, or both. In the conventional formulation this redundancy is not used at
the outer algebraic level, and the working equations therefore retain an
$O(\ndet^2)$ determinant-pair traversal. Here we present an exact
unique-spin-string reformulation of VBSCF determinant algebra in which the
fundamental objects are the sets of unique $\alpha$ and unique $\beta$ spin
strings rather than the full determinant pairs themselves. Each VB structure is
represented by a coefficient matrix on the unique-spin product space, and the
structure-space overlap, one-electron Hamiltonian, same-spin Hamiltonian, and
exact opposite-spin Coulomb term are assembled by matrix contractions on the
compressed spaces. The same reorganization extends to analytic gradients: the
selected-state adjoint is compressed directly to unique-spin matrices and exact
packed-pair channels, which removes the full determinant-pair reverse sweep
from the analytic derivative evaluation. The resulting exact formulation replaces
the determinant-pair bottleneck by contractions governed by $\na$, $\nb$, and
the packed active-pair dimension, while remaining numerically identical to the
reference determinant-pair implementation. This work establishes an exact
matrix-form route to substantially more efficient nonorthogonal VBSCF
determinant algebra.
\end{abstract}

\keywords{valence-bond self-consistent field, nonorthogonal wave functions, determinant algebra, analytic gradients, exact integrals}

\begin{document}

\section{Introduction}
Valence-bond self-consistent field (VBSCF) theory remains one of the most
chemically interpretable \textit{ab initio} descriptions of electronic
structure because the wave function is expanded in terms of localized,
nonorthogonal valence-bond structures. In practical applications the structure
coefficients and localized orbitals are optimized variationally, after which
the wave function can be interpreted directly in terms of resonance,
spin-coupling patterns, and bond rearrangements.

The computational difficulty is equally familiar. Once the orbitals belonging
to different structures are nonorthogonal, the structure-space overlap and
Hamiltonian matrices are usually assembled through expansions over full Slater
determinants. If the $I$th VB structure is written as
\begin{equation}
|\Phi_I\rangle = \sum_d t_{I,d} |D_d\rangle,
\label{eq:structure-expansion}
\end{equation}
then the corresponding structure-space matrix element is
\begin{equation}
M_{IJ} = \sum_{d,e} t_{I,d} M_{de} t_{J,e},
\label{eq:structure-from-determinants}
\end{equation}
where $M_{de}$ denotes the determinant-pair overlap or Hamiltonian kernel. In a
straightforward implementation, Eq.~\eqref{eq:structure-from-determinants}
implies an outer traversal over all full determinant pairs, and the same
determinant-pair structure reappears in analytic derivative expressions. For
large VB expansions this determinant-pair organization becomes the dominant
algebraic bottleneck.

The crucial observation underlying the present work is that the full
determinant list is highly redundant at the spin-resolved level. A full
determinant can be written as
\begin{equation}
D_p = \left( \alpha_{i(p)}, \beta_{j(p)} \right),
\label{eq:full-det-factorization}
\end{equation}
where $\alpha_{i(p)}$ and $\beta_{j(p)}$ are the occupied $\alpha$- and
$\beta$-spin strings. Many different full determinants share identical
$\alpha$-spin strings, identical $\beta$-spin strings, or both. Conventional
VBSCF determinant algebra nevertheless keeps the full determinant pair as the
outer object and therefore revisits identical same-spin subproblems many times.

In this article we show that the outer algebra itself can be reformulated
exactly on the spaces of \emph{unique} $\alpha$-spin strings and \emph{unique}
$\beta$-spin strings. Each structure is represented by a coefficient matrix on
the unique-spin product space, and the overlap, one-electron Hamiltonian,
same-spin Hamiltonian, and exact opposite-spin Coulomb coupling are assembled
directly by matrix contractions on that compressed representation. The same
idea extends to the analytic gradient, where the selected-state adjoint can be
compressed to unique-spin matrices and exact packed-pair channels before the
final backpropagation to orbital and integral derivatives.

The present contribution is therefore an exact algorithmic reformulation rather
than a local cache optimization. The unique-spin-string representation, the
removal of the full determinant-pair outer traversal from the forward
evaluation, and the corresponding matrix-form adjoint equations are all part
of the method introduced here. In the following sections we derive the exact
theory, analyze the resulting scaling, and summarize the numerical behavior of
the implementation with exact active-space two-electron integrals.

\section{Theory}

\subsection{Full determinant formulation}
We begin from the conventional determinant-level formulation. Let the full
determinant expansion contain $\ndet$ determinants and let
\begin{equation}
D_p = \left( \alpha_{i(p)}, \beta_{j(p)} \right)
\end{equation}
denote the decomposition of determinant $p$ into its spin-resolved occupied
strings. The exact determinant-pair overlap factorizes as
\begin{equation}
S_{pq} =
S^{\alpha}_{i(p)i(q)} S^{\beta}_{j(p)j(q)},
\label{eq:det-overlap-factor}
\end{equation}
where $S^{\alpha}$ and $S^{\beta}$ are ordered same-spin determinant overlap
matrices. The one-electron and same-spin Hamiltonian contributions have the
analogous separable form
\begin{equation}
H^{1e}_{pq} =
H^{1e,\alpha}_{i(p)i(q)} S^{\beta}_{j(p)j(q)}
+
S^{\alpha}_{i(p)i(q)} H^{1e,\beta}_{j(p)j(q)},
\label{eq:det-h1-factor}
\end{equation}
and
\begin{equation}
H^{\mathrm{same}}_{pq} =
H^{\alpha}_{i(p)i(q)} S^{\beta}_{j(p)j(q)}
+
S^{\alpha}_{i(p)i(q)} H^{\beta}_{j(p)j(q)}.
\label{eq:det-same-factor}
\end{equation}

The opposite-spin Coulomb coupling is the nontrivial term. In the standard
determinant-pair formulation it can be written exactly in packed active-pair
notation as
\begin{equation}
H^{\alpha\beta}_{pq}
=
\sum_{P=1}^{\npair}
\sum_{Q=1}^{\npair}
U^{\alpha}_{Q}\!\bigl(i(p),i(q)\bigr)
\, G_{QP} \,
U^{\beta}_{P}\!\bigl(j(p),j(q)\bigr),
\label{eq:opposite-spin-standard}
\end{equation}
where $\npair = n_{\mathrm{act}}(n_{\mathrm{act}}+1)/2$ is the number of packed
active orbital pairs, $G_{QP}$ is the exact active-space two-electron kernel in
that packed representation, and $U^{\alpha}_{Q}$ and $U^{\beta}_{P}$ are the
ordered same-spin first-order cofactor projections associated with the
determinant pairs $(\alpha_{i(p)},\alpha_{i(q)})$ and
$(\beta_{j(p)},\beta_{j(q)})$. Equation \eqref{eq:opposite-spin-standard} is
still exact, but when inserted into Eq.~\eqref{eq:structure-from-determinants}
it leaves the full determinant pair as the outer algebraic object.

\subsection{Unique-spin spaces and structure coefficient matrices}
Let $\na$ and $\nb$ denote the numbers of unique $\alpha$-spin and unique
$\beta$-spin occupied strings present in the full determinant expansion. We
define unique-spin index maps
\begin{equation}
p \mapsto a(p) \in \{1,\dots,\na\},
\qquad
p \mapsto b(p) \in \{1,\dots,\nb\},
\label{eq:unique-spin-maps}
\end{equation}
that associate each full determinant with the corresponding unique-spin IDs.

For each VB structure $I$, we then define a coefficient matrix
\begin{equation}
\Cmat_I \in \mathbb{R}^{\na \times \nb},
\end{equation}
with elements
\begin{equation}
\left( \Cmat_I \right)_{ab}
=
\sum_{d:\, a(d)=a,\; b(d)=b} t_{I,d}.
\label{eq:structure-coefficient-matrix}
\end{equation}
Equation \eqref{eq:structure-coefficient-matrix} is simply a regrouping of the
same determinant expansion coefficients that appear in
Eq.~\eqref{eq:structure-expansion}. The only new ingredient is that all full
determinants sharing the same unique-spin pair $(a,b)$ are summed into the same
matrix element of $\Cmat_I$.

This definition immediately replaces the determinant-indexed structure
expansion by a matrix-indexed structure expansion. The crucial point is that
all determinant kernels now depend only on the ordered unique-spin pair IDs,
not on the original full determinant indices. As a result, the structure-space
matrices can be assembled directly from $\{ \Cmat_I \}$ and the ordered
unique-spin channels.

\subsection{Matrix-form structure assembly}
The overlap follows immediately from Eq.~\eqref{eq:det-overlap-factor}. After
regrouping the determinant coefficients according to
Eq.~\eqref{eq:structure-coefficient-matrix}, one obtains
\begin{equation}
S_{IJ}
=
\Frob{\Cmat_I}{\Sa \Cmat_J \Sb^{T}}.
\label{eq:structure-overlap-matrix-form}
\end{equation}
The one-electron contribution is
\begin{equation}
H^{1e}_{IJ}
=
\Frob{\Cmat_I}{\Honea \Cmat_J \Sb^{T}}
+
\Frob{\Cmat_I}{\Sa \Cmat_J \Honeb^{T}},
\label{eq:structure-h1-matrix-form}
\end{equation}
and the same-spin Hamiltonian contribution is
\begin{equation}
H^{\mathrm{same}}_{IJ}
=
\Frob{\Cmat_I}{\Ha \Cmat_J \Sb^{T}}
+
\Frob{\Cmat_I}{\Sa \Cmat_J \Hb^{T}}.
\label{eq:structure-same-matrix-form}
\end{equation}

To treat the exact opposite-spin term, the same-spin first-order cofactor
payloads are regrouped by packed active-pair channel. For each packed pair
$P$, we define the ordered unique-spin matrices
\begin{equation}
\Ua{P} \in \mathbb{R}^{\na \times \na},
\qquad
\mathbf{U}^{\beta}_{P} \in \mathbb{R}^{\nb \times \nb},
\label{eq:packed-pair-channels}
\end{equation}
with matrix elements
\begin{equation}
\left[\Ua{P}\right]_{aa'} = U^\alpha_P(a,a'),
\qquad
\left[\mathbf{U}^{\beta}_{P}\right]_{bb'} = U^\beta_P(b,b').
\label{eq:packed-pair-channel-elements}
\end{equation}
The exact two-electron kernel is then applied at the packed-pair level,
\begin{equation}
\Ub{P} =
\sum_{Q=1}^{\npair}
G_{PQ}\, \mathbf{U}^{\beta}_{Q},
\label{eq:beta-projected-channel}
\end{equation}
so that the opposite-spin structure contribution becomes
\begin{equation}
H^{\alpha\beta}_{IJ}
=
\sum_{P=1}^{\npair}
\Frob{\Cmat_I}{\Ua{P}\, \Cmat_J\, \Ub{P}^{T}}.
\label{eq:structure-opposite-spin-matrix-form}
\end{equation}
Equations \eqref{eq:packed-pair-channels}--\eqref{eq:structure-opposite-spin-matrix-form}
show that the exact opposite-spin term is reorganized into a sum over packed
channels acting on coefficient matrices over the unique-spin spaces, rather
than into a sum over full determinant pairs.

The total structure Hamiltonian matrix is therefore
\begin{equation}
H_{IJ}
=
H^{\mathrm{same}}_{IJ}
+
H^{\alpha\beta}_{IJ}.
\label{eq:structure-total-h}
\end{equation}
Equations \eqref{eq:structure-overlap-matrix-form}--\eqref{eq:structure-total-h}
show that the exact structure-space algebra can be performed directly on the
compressed unique-spin representation. No outer traversal over all full
determinant pairs is required in the exact forward build.

\subsection{Selected-state matrix-form adjoint}
The same compression applies to analytic gradients. Let the state-averaged
target be
\begin{equation}
E_{\mathrm{SA}} = \sum_{n=1}^{\nstate} \wstate \En,
\label{eq:state-average-energy}
\end{equation}
where $\wstate$ are normalized state weights and $\En$ are the generalized
eigenvalues. For each state, the determinant coefficients are gathered into a
state-dependent matrix
\begin{equation}
\Cmat^{(n)} \in \mathbb{R}^{\na \times \nb},
\end{equation}
defined analogously to Eq.~\eqref{eq:structure-coefficient-matrix} but using
the determinant coefficients of the selected state rather than the
structure-expansion coefficients.

Because the generalized eigenvalue derivative carries coefficients $\wstate$
for Hamiltonian terms and $-\wstate \En$ for overlap terms, the same-spin and
one-electron adjoints compress to unique-spin weight matrices
\begin{equation}
\mathbf{W}^{(H)}_{\alpha}
=
\sum_n \wstate \,
\Cmat^{(n)} \Sb [\Cmat^{(n)}]^{T},
\label{eq:alpha-H-weight}
\end{equation}
\begin{equation}
\mathbf{W}^{(S)}_{\alpha}
=
\sum_n (-\wstate \En)\,
\Cmat^{(n)} \Sb [\Cmat^{(n)}]^{T},
\label{eq:alpha-S-weight}
\end{equation}
with the $\beta$-spin counterparts
\begin{equation}
\mathbf{W}^{(H)}_{\beta}
=
\sum_n \wstate \,
[\Cmat^{(n)}]^{T} \Sa \Cmat^{(n)},
\label{eq:beta-H-weight}
\end{equation}
\begin{equation}
\mathbf{W}^{(S)}_{\beta}
=
\sum_n (-\wstate \En)\,
[\Cmat^{(n)}]^{T} \Sa \Cmat^{(n)}.
\label{eq:beta-S-weight}
\end{equation}
and analogous partner-transfer matrices are obtained by replacing $\Sa$ or
$\Sb$ by the appropriate same-spin Hamiltonian blocks. The reverse
accumulation to orbital-overlap, one-electron, and same-spin two-electron
derivatives is therefore carried out over ordered unique-spin pairs rather than
over full determinant pairs.

The opposite-spin adjoint is similarly matrix-form. If
$\mathbf{U}^{\alpha}_{P}$ and $\mathbf{U}^{\beta}_{Q}$ denote the packed-pair
channel families regrouped from the same-spin caches, then the adjoint with
respect to the packed active-space two-electron kernel can be written as
\begin{equation}
\frac{\partial E}{\partial G_{QP}}
=
\sum_n \wstate\,
\Frob{
\mathbf{U}^{\alpha}_{P}
}{
\Cmat^{(n)} \mathbf{U}^{\beta}_{Q} [\Cmat^{(n)}]^{T}
}.
\label{eq:opposite-spin-adjoint}
\end{equation}
The overlap adjoint on each spin sector follows from the same channel family
after replacing first-order cofactor projections by inverse-overlap projected
channels. Thus, the reverse sweep is reorganized on the same unique-spin
objects that appear in the exact structure build.

\subsection{Complexity and memory analysis}
The original determinant-pair assembly retains an outer complexity skeleton
of
\begin{equation}
O(\ndet^2),
\label{eq:old-scaling}
\end{equation}
because the unreduced workflow is organized around full determinant pairs. In
contrast, the present reformulation separates the exact workflow into three
parts:
\begin{enumerate}
  \item construction of ordered unique same-spin channels,
  \item matrix-form structure assembly or selected-state compression on the
        $\na \times \nb$ coefficient matrices, and
  \item final backpropagation over ordered unique-spin pairs.
\end{enumerate}

The ordered same-spin cache scales as
\begin{equation}
O(\na^2 + \nb^2),
\label{eq:cache-scaling}
\end{equation}
whereas the structure-space matrix assembly is controlled by matrix products of
the form $\mathbf{A}\Cmat_I\mathbf{B}^{T}$:
\begin{equation}
O\!\left(
\nstr \left( \na^2 \nb + \na \nb^2 \right)
\right),
\label{eq:forward-scaling}
\end{equation}
plus the packed-pair channel sum for the opposite-spin contribution,
\begin{equation}
O\!\left(
\nstr \npair \left( \na^2 \nb + \na \nb^2 \right)
\right).
\label{eq:forward-op-scaling}
\end{equation}
For the selected-state adjoint, the same matrix products appear with $\nstate$
in place of $\nstr$:
\begin{equation}
O\!\left(
\nstate \npair \left( \na^2 \nb + \na \nb^2 \right)
\right),
\label{eq:backward-op-scaling}
\end{equation}
followed by scatter operations over ordered unique-spin pairs. Thus, the
forward and backward equations are no longer controlled by an explicit
$O(\ndet^2)$ outer traversal.

The memory consequences are also important. A naive scalar opposite-spin tensor
defined directly on the full unique-spin pair grid would have size
\begin{equation}
O(\na^2 \nb^2),
\label{eq:old-op-cache}
\end{equation}
which is unnecessary once the forward and backward equations consume the
per-spin projected packed-pair channels directly. In the present
implementation, the matrix-form algorithm stores the ordered same-spin
pair caches and the projected packed-pair channel families, while the old
scalar opposite-spin cache is optional and retained only for legacy pairwise
diagnostics.

\section{Implementation}
The method was implemented in the \textsc{xmvb-cpp} VBSCF code as a new exact
formulation for both the forward structure-space matrix construction and the
analytic active-space adjoint. The implementation is organized around a
unique-spin-string context that stores
\begin{enumerate}
  \item the unique-spin index maps from full determinants to unique $\alpha$-
        and
        $\beta$-spin strings,
  \item ordered unique same-spin pair evaluations, and
  \item exact packed-pair projections needed by the opposite-spin channels.
\end{enumerate}

The structure-space forward build proceeds as follows.
\begin{enumerate}
  \item The determinant expansion is compressed into one coefficient matrix
        $\Cmat_I$ per structure.
  \item The ordered same-spin caches are reorganized into dense overlap and
        Hamiltonian matrices on the unique-spin spaces.
  \item The opposite-spin first-order cofactor payload is regrouped by packed
        active-pair channel.
  \item The structure-space overlap and Hamiltonian matrices are assembled by
        matrix products and Frobenius inner products, following
        Eqs.~\eqref{eq:structure-overlap-matrix-form}--\eqref{eq:structure-total-h}.
\end{enumerate}

The analytic gradient follows an analogous sequence.
\begin{enumerate}
  \item The selected-state determinant coefficients are compressed into
        state-specific matrices $\Cmat^{(n)}$.
  \item The same-spin / one-electron and opposite-spin adjoints are evaluated
        by matrix-form contractions on the unique-spin spaces.
  \item The resulting weight matrices are backpropagated through the ordered
        same-spin caches to obtain the overlap, one-electron, and active-space
        two-electron gradients.
\end{enumerate}

This design preserves exactness while moving the determinant algebra to the
relevant compressed spaces. The legacy pairwise build is still retained for
diagnostic interfaces that explicitly request storage of each full
determinant-pair evaluation. Thus, the new formulation augments rather than
removes the reference path.

\section{Computational Details}
All calculations were carried out with the \textsc{xmvb-cpp} implementation of
the present exact unique-spin VBSCF algorithm. Exact active-space two-electron
integrals were generated from the \textsc{libcint} integral engine. The
generalized eigenvalue problem for the structure-space overlap and Hamiltonian
matrices was solved at each SCF step in the usual way. Approximate low-rank
variants are left for future work.

\todo{Insert exact hardware information, compiler version, BLAS/Eigen setup,
and thread-count policy.}

\todo{Insert basis sets, active spaces, and structure-selection details for the
final benchmark set.}

\section{Results and Discussion}

\subsection{Algebraic verification}
The first test of the reformulation is strict agreement with the reference full
determinant-pair implementation. For a representative exact-integral benchmark,
the matrix-form structure build reproduces the reference overlap, one-electron,
and total Hamiltonian matrices to machine precision. In one representative
benchmark, the determinant space contains $\ndet = 400$ full determinants but
only $\na = 20$ unique
$\alpha$-spin strings and $\nb = 20$ unique $\beta$-spin strings. In that
case, the ordered same-spin cache requires only
\begin{equation}
\na^2 + \nb^2 = 20^2 + 20^2 = 800
\end{equation}
ordered same-spin evaluations, compared with
\begin{equation}
2 \times \frac{\ndet(\ndet+1)}{2} = 160{,}400
\end{equation}
same-spin determinant kernels in the unreduced full determinant-pair picture.

\todo{Insert a concise table comparing matrix-form and reference structure
matrices. Current internal checks show machine-precision agreement.}

\subsection{Analytic gradient verification}
The matrix-form adjoint was validated by finite-difference comparisons for both
the active-space two-electron and orbital gradients. In the current exact-mode
tests on \ce{F2}, the same-spin and opposite-spin matrix-form adjoints recover
the numerical derivatives with errors on the order of
\num{1e-8}--\num{1e-9}, confirming that the compressed reverse-mode algebra is
exact in the regime studied here.

\todo{Replace this paragraph by a final table of finite-difference checks for
all benchmark systems in the exact-integral regime.}

\subsection{Benchmark behavior}
The most important performance result is that the same-spin determinant algebra
is no longer the leading cost of a VBSCF step. In representative exact
calculations, same-spin cache construction, structure-space matrix build, and
the matrix-form adjoint become small relative to the remaining AO- and
active-space integral kernels. For the $\ndet=400$, $\na=\nb=20$ example
above, the exact structure build is already subdominant to the
construction of active-space two-electron intermediates and to AO-level
backpropagation.

\begin{table}[t]
  \caption{Suggested benchmark summary table for the final manuscript.
  Replace the placeholder entries by the final measured values for the exact
  implementation.}
  \label{tab:benchmark-summary}
  \begin{tabular}{lccccc}
    \toprule
    System & $\ndet$ & $\na$ & $\nb$ & Stage & Time (s) \\
    \midrule
    \todo{Case A} & \todo{...} & \todo{...} & \todo{...} & Structure build & \todo{...} \\
    \todo{Case A} & \todo{...} & \todo{...} & \todo{...} & Adjoint & \todo{...} \\
    \todo{Case B} & \todo{...} & \todo{...} & \todo{...} & Structure build & \todo{...} \\
    \todo{Case B} & \todo{...} & \todo{...} & \todo{...} & Adjoint & \todo{...} \\
    \bottomrule
  \end{tabular}
\end{table}

\begin{figure}[t]
  \centering
  \fbox{\rule{0pt}{2.0in}\rule{0.9\linewidth}{0pt}}
  \caption{Suggested schematic figure. A useful final figure would show the
  mapping from full determinants $D_d$ to the unique-spin index maps $a(d)$ and
  $b(d)$,
  the resulting structure coefficient matrices $\Cmat_I$, and the matrix-form
  contractions used for the overlap, same-spin, and opposite-spin channels.}
  \label{fig:scheme}
\end{figure}

\begin{figure}[t]
  \centering
  \fbox{\rule{0pt}{2.0in}\rule{0.9\linewidth}{0pt}}
  \caption{Suggested scaling plot. Compare the wall time of the reference
  determinant-pair build and the present matrix-form build as a function of
  $\ndet$ and/or the compression ratios
  $\ndet/\na$ and $\ndet/\nb$.}
  \label{fig:scaling}
\end{figure}

\subsection{Relation to broader nonorthogonal methods}
The present reformulation is complementary to the broader development of
nonorthogonal electronic structure theory. Generalized nonorthogonal matrix
element formalisms seek compact ways to couple arbitrary nonorthogonal
determinants and excitations. Our emphasis here is different: rather than
generalizing the determinant-pair kernel itself, we exploit the special product
structure of the VBSCF determinant expansion to reorganize the many repeated determinant-pair
contractions into matrix algebra on the unique-spin spaces. The resulting
algorithm is therefore exact for the VBSCF determinant algebra considered here,
yet conceptually compatible with future extensions that combine unique-spin
compression with more general nonorthogonal matrix-element technology.

\section{Discussion}
From a methodological perspective, the most important feature of the present
approach is that it is not a local implementation tweak layered on top of the
original determinant-pair algorithm. The unique-spin-string compression, the
matrix-form structure assembly, and the matrix-form backward equations are new
components introduced together here, so that the \emph{outer algebraic object}
is no longer the full determinant pair. This change of algebraic organization
is the essential source of the removal of the $O(\ndet^2)$
determinant-pair bottleneck.

The implementation also clarifies where further work is most valuable. Once the
same-spin and opposite-spin determinant algebra is compressed to the unique-spin
spaces, the dominant cost in many exact calculations shifts to the construction
of active-space two-electron intermediates and to AO-level backpropagation.
Future efforts can therefore focus on these remaining kernels without
revisiting the determinant-algebra bottleneck.

Two natural extensions deserve special mention. First, sparse or block-sparse
representations of the structure coefficient matrices $\Cmat_I$ may offer
further speedups when the determinant expansion occupies only a small subset of
the unique-spin product grid. Second, the present exact framework provides a
natural starting point for future low-rank treatments of the opposite-spin
kernel, including factorized or Cholesky-compressed channel representations. These
approximations are not part of the present work; here the emphasis is the exact
unique-spin reformulation itself. Third, the same matrix-form machinery should
also be useful for selected-determinant, adaptive-structure, and related
NOCI-style extensions in which the determinant algebra remains nonorthogonal
but exhibits strong internal shared-spin structure.

\section{Conclusions}
We have presented an exact unique-spin reformulation of determinant algebra in
VBSCF. The key step is to replace the full determinant-pair outer loop by a
matrix representation on the spaces of unique $\alpha$- and $\beta$-spin
strings. Each structure and selected state is represented by a coefficient
matrix on the unique-spin product space, and the structure-space overlap,
one-electron Hamiltonian, same-spin Hamiltonian, exact opposite-spin Coulomb
term, and their analytic adjoints are assembled by matrix contractions on that
compressed representation.

This reformulation removes the determinant-pair bottleneck from the forward and
backward evaluations while remaining exact with respect to the reference
active-space integral formulation. Numerical tests confirm machine-precision
agreement with the reference determinant-pair build and finite-difference
agreement for the analytic gradients. In practical applications, the
same-spin/opposite-spin determinant algebra ceases to be the leading cost of a
VBSCF step, exposing the remaining AO- and active-space integral kernels as the
next targets for optimization.

More broadly, the present work demonstrates that nonorthogonal VBSCF
determinant algebra admits an exact matrix-form organization that is much
closer to the true shared-spin structure of the method than the traditional full
determinant-pair formulation. We expect this viewpoint to be useful not only
for high-performance VBSCF implementations but also for future developments in
nonorthogonal multireference and NOCI-style electronic structure methods. In
future work, it should also provide a clean exact foundation for low-rank
compression strategies in the opposite-spin sector.

\section*{Supporting Information}
\begin{itemize}
  \item Additional derivations for the opposite-spin packed-pair channels.
  \item Detailed finite-difference validation tables.
  \item Additional timing decompositions and scaling plots.
  \item Input files and benchmark settings for the reported test cases.
\end{itemize}

\section*{Author Information}
\subsection*{Corresponding Author}
\begin{itemize}
  \item \todo{Name, e-mail, ORCID.}
\end{itemize}

\subsection*{Author Contributions}
\todo{Write the final author-contribution statement.}

\subsection*{Notes}
The authors declare no competing financial interest.

\section*{Acknowledgement}
\todo{Insert funding, computing resources, and project acknowledgements.}

\begin{tocentry}
\centering
\fbox{\rule{0pt}{1.6in}\rule{0.9\linewidth}{0pt}}

\vspace{0.5em}
\small
Unique-spin coefficient matrices reorganize the exact VBSCF determinant algebra
from full determinant pairs to matrix contractions on the spaces of unique
$\alpha$ and $\beta$ spin strings.
\end{tocentry}

\end{document}
